
Let \(X\) be a formal variety over \(k\). If \(x\in X\) projects to
\(s\in\operatorname{Spf}(k)\), then the residue field \(\kappa(x)\)
is a finite extension of \(\kappa(s)\). We denote by
\(\kappa_e(x)\) the separable closure of \(\kappa(s)\) in
\(\kappa(x)\).\footnote{\label{sga3:VIIB:nde:73}In what follows, the
editors used the preceding additions to give details of the
construction of the functor \(X\mapsto X_e\) and to prove that it
commutes with finite products; this is used in \SGARefLink{sga3:VIIB:2.5.1:corollary}{2.5.1}.}

\medskip
\noindent\textbf{Proposition 1.6.F.}
\phantomsection
\label{sga3:VIIB:1.6.F}
{\itshape \SGARefLink{sga3:vii:paragraph:5-1:equation:i}{(i)} The inclusion
\[
  \operatorname{Vaf}^{\mathrm{\acute{e}t}}/k
  \longrightarrow
  \operatorname{Vaf}/k
\]
has a left adjoint \(X\mapsto X_e\). The formal variety \(X_e\) has
the same points as \(X\). For each \(x\in X\) projecting to
\(s\in\operatorname{Spf}(k)\), choose a primitive element
\(\xi\) of \(\kappa_e(x)/\kappa(s)\), let \(n\) be its degree, choose
any lift \(u\in\mathcal O_{X,x}\) of \(\xi\), and let
\(F\in k[T]\) be a monic polynomial of degree \(n\) that annihilates
\(u\). Then put
\[
  \mathcal O_{X_e,x}=k[T]/(F).
\]
For every
\(Y\in\operatorname{Ob}(\operatorname{Vaf}^{\mathrm{\acute{e}t}}/k)\),
the canonical maps in the following diagram are bijective:
\begin{center}
  \phantomsection
  \label{sga3:VIIB:1.6.F:hom-square}
  \begin{tikzcd}[row sep=2.8em,column sep=3.5em]
\operatorname{Hom}_{\operatorname{Vaf}/k}(X,Y)
  \arrow[r,"{\sim}"]
  & \operatorname{Hom}_{\operatorname{Vaf}/\kappa(k)}
      (X_\kappa,Y_\kappa)
      \arrow[d,"\wr"] \\
\operatorname{Hom}_{\operatorname{Vaf}/k}(X_e,Y)
  \arrow[r,"{\sim}"]
  & \operatorname{Hom}_{\operatorname{Vaf}/\kappa(k)}
      ((X_e)_\kappa,Y_\kappa)\, ,
\end{tikzcd}

\end{center}
The vertical arrow is induced by the inclusions
\[
  \kappa_e(x)\lhook\joinrel\longrightarrow\kappa(x),
  \qquad x\in X.
\]
In particular, this defines a morphism
\[
  p:X\longrightarrow X_e,
\]
and every morphism of formal \(k\)-varieties
\(\phi:X\to Y\), with \(Y\) étale over \(k\), factors uniquely
through \(p\).

\SGARefLink{sga3:vii:paragraph:5-1:equation:ii}{(ii)} The functor \(X\mapsto X_e\) commutes with finite products.}

\medskip
\noindent\textit{Proof.}
Assertion~\textup{\SGARefLink{sga3:vii:paragraph:5-1:equation:i}{(i)}} follows from \SGARefLink{sga3:VIIB:1.6.E}{1.6.E}~\textup{(b)} and \SGARefLink{sga3:VIIB:1.6.D}{1.6.D}.
We prove~\textup{\SGARefLink{sga3:vii:paragraph:5-1:equation:ii}{(ii)}}. By the equivalence of categories in
\SGARefLink{sga3:VIIB:1.6.E}{1.6.E}~\textup{(b)}, we may suppose that \(k\) is a field. In that
case, if \(X\) is a semilocal formal \(k\)-variety, so that its
affine algebra \(A\) is semilocal, the affine algebra of \(X_e\) is
the largest subalgebra of \(A\) that is étale over \(k\); denote it
by \(A_e\).

It is therefore enough to show that, if \(K\) and \(L\) are two
finite extensions of \(k\), then
\[
  K_e\otimes_kL_e
  =
  (K\otimes_kL)_e.
\]
Let \(p\) be the characteristic exponent of \(k\), so that \(p=1\)
if \(\operatorname{char}(k)=0\), and
\(p=\operatorname{char}(k)\) otherwise. For every
\(x\in K\otimes_kL\), there is an integer \(n\geq1\) such that
\[
  x^{p^n}\in K_e\otimes_kL_e.
\]
Thus \(K\otimes_kL\) is radicial over \(K_e\otimes_kL_e\), and its
largest étale \(k\)-subalgebra is precisely
\(K_e\otimes_kL_e\). Hence
\[
  K_e\otimes_kL_e=(K\otimes_kL)_e.
\]

Keeping this notation, observe that
\(\mathcal O_{X_e,x}\) is not necessarily a subalgebra of
\(\mathcal O_{X,x}\). It is a subalgebra when \(X\) is
topologically flat over \(k\), by the proposition that follows.
